\section{ANTECEDENTES Y ESTADO DEL ARTE}
\label{sec:antecedentes}

La revisión del estado del arte para el control y la gestión ciber-física de microrredes eléctricas aisladas se articula en torno a cinco ejes fundamentales: la evolución del control secundario por consenso en Sistemas Multiagente (MAS), el despliegue de inteligencia en el borde (\textit{Edge AI/IoT}), el modelado de la resiliencia telemétrica ante retardos y fallas en la comunicación, el paradigma de co-simulación ciber-física multitasa y la integración de dinámicas heterogéneas de fuentes renovables.

\subsection{Evolución del Control Secundario Distribuido y Consenso en Sistemas Multiagente}
La transición desde esquemas de control centralizado hacia arquitecturas distribuidas ha posicionado a la teoría de consenso sobre grafos algebraicos como la metodología predominante para la restauración de voltaje y frecuencia en microrredes \citep{Mamun2026, Altin2023MultiAgent}. El origen de esta línea se sitúa en las leyes de promedio distribuido que restablecen frecuencia y tensión en microrredes isleñas sin requerir un controlador central \citep{SimpsonPorco2015}, y su desarrollo posterior se ordena habitualmente en tres familias ---centralizada, distribuida y descentralizada--- según el grado de dependencia de un nodo coordinador y el alcance de la información compartida \citep{Khayat2020, Li2019}. En el enfoque clásico los agentes actualizan sus variables de control de forma asintótica en un horizonte de tiempo infinito. La búsqueda de garantías temporales explícitas condujo primero a la estabilidad en tiempo finito, cuyas cotas dependen del estado inicial, y luego a la estabilidad en tiempo fijo, cuya base teórica es el diseño de realimentación no lineal con exponentes fraccionarios duales que acota la convergencia mediante un $T_{\max}$ independiente de las condiciones iniciales \citep{Polyakov2012}. Sobre esa base se consolidaron protocolos de consenso en tiempo finito para frecuencia y tensión que incorporan retardos de comunicación \citep{Ning2021} y mecanismos de disparo por eventos \citep{Choi2022}. La frontera actual la constituyen los esquemas en tiempo predefinido, donde el instante de convergencia es un parámetro de diseño explícito \citep{Huang2023}, y los protocolos de convergencia en tiempo exacto \citep{Amrr2026}; en esa familia se inscribe el antecedente más cercano al protocolo aquí empleado \citep{Shrivastava2023FiniteTime}. Conviene precisar que el aporte de este trabajo no consiste en una nueva cota temporal, sino en la integración de estos protocolos con una plataforma de co-simulación multitasa desplegada sobre hardware de borde. En el eje específico de regulación de tensión y reparto de potencia reactiva, la literatura reporta esquemas resilientes ante perturbaciones y fallas de comunicación \citep{Ge2021}, estrategias de regulación de barra crítica en microrredes múltiples con tiempos mínimos positivos entre eventos \citep{Zhai2024} y esquemas distribuidos para topologías multibarrables con reparto adaptativo de reactiva \citep{Huang2026TIE}.

Paralelamente, para optimizar el esfuerzo de control en sistemas con dinámicas heterogéneas, se han propuesto formulaciones unificadas de consenso óptimo que combinan despacho económico con esquemas de seguimiento de líder (\textit{leader-follower}) y arquitecturas libres de líder (\textit{leaderless}), permitiendo a los convertidores modular sus trayectorias ante variaciones operativas abruptas \citep{Guo2026, Salman2024}. Asimismo, la literatura reciente explora enfoques adaptativos libres de modelo (\textit{Model-Free Adaptive Control}, MFAC) y de regulación cooperativa basados exclusivamente en datos (\textit{Data-Driven Control}), los cuales permiten diseñar leyes de control tolerantes a fallos sin requerir la parametrización exacta ni la topología física completa de la microrred \citep{Zhao2025, Xu2025, Ren2024DataDriven}.

Para resolver la sobrecarga telemétrica y la congestión de los canales digitales, diversos autores han incorporado mecanismos de consenso activados por eventos (\textit{event-triggered consensus}) \citep{Chen2024}. En estas formulaciones, los agentes transmiten sus estados telemétricos únicamente cuando la desviación respecto a su último estado transmitido supera un umbral dinámico de Lyapunov, reduciendo drásticamente el tráfico de red mientras se preserva la estabilidad asintótica de la frecuencia y el voltaje nominal \citep{Liu2024, AlTameemi2024}.

La tabla~\ref{tab:comparacion_control} sintetiza las arquitecturas revisadas frente a los requisitos que impone esta tesis.

\begin{table}[H]
    \centering
    \footnotesize
    \begin{tabularx}{\textwidth}{|>{\raggedright\arraybackslash}p{3.0cm}|>{\centering\arraybackslash}X|>{\centering\arraybackslash}X|>{\centering\arraybackslash}X|>{\centering\arraybackslash}X|>{\centering\arraybackslash}X|}
        \hline
        \textbf{Arquitectura} & \textbf{Coordinador central} & \textbf{Instante de convergencia} & \textbf{Cota independiente del estado inicial} & \textbf{Considera retardos} & \textbf{Pondera la antigüedad del dato} \\
        \hline
        Centralizada (OPF/MPC) \citep{Khayat2020, Guo2026} & Sí & No & No & No & No \\
        \hline
        Líder-seguidor \citep{Guo2026} & Sí & Según diseño & No & Parcial & No \\
        \hline
        Estatismo (\textit{droop}) \citep{Khayat2020} & No & No & No & No aplica & No \\
        \hline
        Consenso asintótico \citep{SimpsonPorco2015} & No & No & No & Parcial & No \\
        \hline
        Consenso en tiempo finito \citep{Ning2021, Shrivastava2023FiniteTime} & No & Sí & No & Sí & No \\
        \hline
        \textbf{Consenso en tiempo fijo (este trabajo) \citep{Polyakov2012, Huang2023}} & \textbf{No} & \textbf{Sí} & \textbf{Sí} & \textbf{Sí} & \textbf{Sí} \\
        \hline
    \end{tabularx}
    \caption{Comparación de las arquitecturas de control secundario frente a los requisitos de este trabajo. La última fila corresponde al esquema desarrollado en esta tesis.}
    \label{tab:comparacion_control}
\end{table}

\subsection{Inteligencia en el Borde y Arquitecturas Ciber-Físicas IoT}
El despliegue de Sistemas Multiagente en sistemas de potencia ha evolucionado de simulaciones centralizadas de software hacia arquitecturas físicas distribuidas en la periferia de la red (\textit{Edge Computing} e \textit{IoT}) \citep{Diaz2025, AlvarezLopez2026, Mehmood2021}. En estos entornos, microcomputadores embebidos de bajo consumo (tales como dispositivos Raspberry Pi o microcontroladores industriales) actúan como nodos inteligentes desacoplados, alojando instancias individuales de los agentes de control que interactúan mediante protocolos de mensajería estandarizados.

La literatura destaca el uso de infraestructuras IoT y plataformas telemétricas ligeras como MQTT y WebSockets para la coordinación descentralizada y el monitoreo en tiempo real de microrredes sostenibles \citep{Diaz2025}. No obstante, la arquitectura dominante en los marcos telemétricos comerciales sigue apoyándose en corredores de mensajería (\textit{brokers}), lo que impone el desafío metodológico de preservar la autonomía de decisión de par a par (\textit{peer-to-peer}, P2P) en la capa algorítmica mientras se aprovecha la interoperabilidad de la capa telemétrica, según documentan \citep{Diaz2025, Mamun2026}.

\subsection{Resiliencia ante Retardos, Pérdidas de Conectividad e Incertidumbre Telemétrica}
Uno de los supuestos más limitantes en la literatura teórica de consenso es la asunción de canales de comunicación ideales, libres de ruido, sin pérdida de paquetes y con latencias despreciables. En sistemas ciber-físicos reales, los retardos de propagación y la pérdida estocástica de información son inherentes a los enlaces inalámbricos o redes cableadas compartidas \citep{Olayanju2025ConsensusFreq}.

Investigaciones experimentales y analíticas demuestran que, mientras los modelos continuos ideales predicen estabilidad ante retardos teóricos de varios segundos, el acoplamiento no lineal con los inversores y las líneas de distribución puede inducir oscilaciones de potencia e inestabilidad armónica ante latencias superiores a 300 ms \citep{Olayanju2025ConsensusFreq}. Para contrarrestar esta vulnerabilidad, se han desarrollado observadores distribuidos por modos deslizantes (\textit{Sliding Mode Observers}, SMO) y bancos de estimación robustos capaces de mitigar ciberataques coordinados de inyección de datos falsos (\textit{False Data Injection Attacks}, FDIA) y aislar enlaces desconectados antes de que provoquen el colapso del consenso \citep{Barzegari2022Resilient, TaverasCruz2025, Fathima2024}. No obstante, la mayoría de estas estrategias reacciona ante el evento declarado de pérdida de paquete o de enlace caído, y no ante el deterioro silencioso de la utilidad temporal del dato recibido.

\noindent Un aspecto menos atendido en la literatura es el tratamiento del envejecimiento de la información de estado. Los esquemas de consenso tolerantes a impedimentos de comunicación suelen modelar retardos acotados y pérdidas estocásticas de paquetes \citep{Olayanju2025ConsensusFreq}, o bien reducir la tasa de transmisión mediante disparadores por eventos que preservan la estabilidad asintótica \citep{Qiu2025, Chen2024}. Sin embargo, estos enfoques asumen implícitamente que todo mensaje recibido sigue siendo representativo del estado actual de su emisor: un paquete retrasado se incorpora al consenso con el mismo peso que uno fresco, de modo que una ráfaga de datos obsoletos puede arrastrar la convergencia hacia un punto de operación inválido sin que el algoritmo lo detecte. La literatura de estabilidad transitoria ha mostrado, desde la teoría de redes complejas y la percolación, que la degradación topológica progresiva ante la pérdida de enlaces compromete la propagación de la referencia antes de que la conectividad se pierda por completo \citep{Machuca2017}. Este trabajo incorpora esa dimensión temporal ausente mediante un mecanismo de desvanecimiento dinámico de enlaces ponderado por la edad del último estado válido recibido, que se detalla en la sección~\ref{sec:implementacion}.

\subsection{Metodologías de Co-Simulación Ciber-Física y Desafíos Multitasa}
Para validar esquemas de control distribuido sin incurrir en simplificaciones que invaliden el anclaje físico de la red (\textit{no physical grounding}), la comunidad científica ha recurrido al paradigma de la \textbf{Co-Simulación Ciber-Física} \citep{Gomes2022FMI30}. En esta metodología, se acoplan resolvedores especializados en diferentes dominios mediante interfaces estándar (como FMI 3.0, sockets TCP/IP o ZeroMQ), permitiendo que un simulador de transitorios electromagnéticos interactúe en tiempo real o cuasi-real con un entorno distribuido de agentes \citep{Muszynski2026EnergyTwin}.

El principal desafío en la co-simulación ciber-física radica en el \textit{desajuste multirate} (\textit{multi-rate mismatch}) derivado de la coexistencia de constantes de tiempo rápidas en la electrónica de potencia (pasos de microsegundos a milisegundos) y tasas de decisión más lentas en el control secundario y despacho (cientos de milisegundos a segundos) \citep{Liu2025Multirate}. La literatura reciente aborda este desacoplamiento mediante esquemas de sincronización con control de paso variable, correctores basados en modelos y retenedores de orden cero (ZOH) o primer orden (FOH), demostrando que una formulación numérica rigurosa es indispensable para evitar la amplificación artificial de errores de discretización en la interfaz de acoplamiento \citep{Haid2022CosimCorrector, Kyllingstad2025CosimErrors}.

\subsection{Modelado de Dinámicas Heterogéneas y Sincronización de Recursos Distribuidos}
La formulación de un banco de pruebas representativo exige incorporar la respuesta dinámica no lineal de los diversos recursos energéticos distribuidos (DER) que componen una microrred aislada:
\begin{itemize}
    \item \textbf{Recursos basados en inversores (IBRs):} Sistemas solares fotovoltaicos y bancos de baterías (BESS), caracterizados por una inercia física nula, dinámicas de transitorios electromagnéticos (EMT) ultrarrápidas y lazos de control en coordenadas síncronas $dq$ \citep{Mehrzad2025, Liu2025Multirate}.
    \item \textbf{Sistemas de sincronización de fase (SRF-PLL):} Implementación de lazos de enganche de fase robustos para garantizar la extracción instantánea de la frecuencia y el ángulo de red bajo condiciones de distorsión armónica o desbalance de carga \citep{Kundu2024ICCFPLL}.
    \item \textbf{Coordinación híbrida AC/DC y generación síncrona de respaldo:} Integración de generadores diésel gobernados que suministran inercia rotacional de base y actúan como nodo de referencia (\textit{slack}) durante la transición en isla, junto con convertidores de interconexión para la gestión óptima de flujos de potencia activa y reactiva \citep{Alam2024, MartinezGomez2025Distributed, Pompodakis2022Sensitivity, DiFazio2024Linear}.
\end{itemize}

\subsection{Identificación de la Brecha y Posicionamiento de la Investigación}
A partir del análisis exhaustivo del estado del arte, se identifican las siguientes limitaciones estructurales en los trabajos previos:
\begin{enumerate}[1.]
    \item \textit{Predominio de simulaciones abstractas puras:} La mayoría de los algoritmos de consenso se evalúan en entornos de software matemático (MATLAB/Simulink aislado o Python puro), omitiendo el despliegue en hardware embebido real y los efectos estocásticos de redes de comunicación compartidas.
    \item \textit{Desconexión entre la capa telemétrica y la física:} Los marcos experimentales con hardware en el bucle rara vez combinan una co-simulación multitasa rigurosa con sincronización determinista entre micro-pasos de inversores ($1\text{ ms}$) y macro-pasos de red ($500\text{ ms}$).
    \item \textit{Falta de evaluación integral ante fallas ciber-físicas combinadas:} Escasean los estudios que sometan simultáneamente al control secundario distribuido a contingencias eléctricas (pérdida súbita de generación renovable, variaciones de irradiancia) y contingencias telemétricas severas (latencia variable, pérdida de paquetes y partición temporal del grafo de comunicación).
\end{enumerate}

La presente tesis de grado se posiciona directamente en esta brecha científica, desarrollando una plataforma integral de control secundario distribuido en tiempo fijo validada sobre un clúster de hardware de borde Raspberry Pi 5 y un entorno de co-simulación ciber-física acoplada, asegurando la resiliencia operativa y el anclaje físico indispensable para su futura transferencia tecnológica.